\clearpage
\sscp{Innovations among the consonants}{08-01-02}
Turning to the consonants,
it was noted above that PMP *y = \it{y}; *z > \it{y}; *j > \it{y, l}, {\0}.
\trf{08-08} shows the reflexes of *y/*z > \it{y} in \tsc{Taliabu}.
Note that a similar merger of *y/*z > **y is found in \tsc{Nuclear
Sula-Buru}, but not Ambelau.
For comparison, in \tsc{Saluan-Banggai} languages, *z > \it{d,
s, ʤ, l} in different languages.
See also reflexes for *qazay in \trf{08-04}.
An additional example of *z > \it{y} in Taliabu is *tuzuk > \it{ka-tuyuk}
‘point at, point out’.

\begin{table}\tcp{\tsc{Taliabu} and nearby languages *y and *z}{08-08}
%\begin{adjustbox}{width=\textwidth}
	\begin{threeparttable}\st{3.5pt}
	\begin{tabular}{lllllll}\lsptoprule
local gl.	&	sail	&	land(ward)	&	big	&	crocodile	&	far	&	rain\\
PMP	&	*la\tbr{y}aR	&	*da\tbr{y}a	&	*Ra\tbr{y}a	&	*buqa\tbr{y}a	&	*\tbr{z}auq	&	*qu\tbr{z}an\\ \midrule
Saluan	&		&		&	\cg{(dakaʔ)}	&	{\hr}buaaʔ	&	{\hr}ma-\tbb{ʤ}oʔon	&	{\hr}u\tbb{ʤ}an\\
Balantak	&		&		&	\cg{(moolaʔ)}	&	{\hr}bua\tbr{y}a	&	{\hr}o-\tbb{l}oa	&	{\hr}u\tbb{s}an\\
Banggai	&	\cg{(kadut)}	&		&		&	{\hr}buea	&	{\hr}o-\tbb{d}oon	&	{\hr}u\tbb{d}an\\
\dshmidrule{7}
Soboyo	&		&	{\hr}ho\tbr{y}o	&	{\hr}ha\tbr{y}a	&	{\hr}fo\tbr{y}a	&	{\hr}\tbr{y}awo	&	{\hr}u\tbr{y}ain\\
Taliabu	&	{\hr}la\tbr{y}a	&	{\hr}ho\tbr{y}o	&	{\hr}ha\tbr{y}a	&	{\hr}fo\tbr{y}a	&	{\hr}\tbr{y}auw	&	{\hr}u\tbr{y}aɲ\\
Kadai	&		&		&	{\hr}ha\tbr{y}a	&	{\hr}fo\tbr{y}a	&	{\hr}\tbr{y}awo	&	{\hr}u\tbr{y}a\\	\midrule
Mangole	&		&	{\hr}lai	&		&	{\hr}fo\tbr{y}a 	&	{\hr}\tbr{y}au	&	{\hr}u\tbr{y}a\\
Sanana	&		&	{\hr}hai	&	{\hr}a\tbr{y}a	&	{\hr}fua\tbr{y}a	&	{\hr}\tbr{y}au	&	{\hr}u\tbr{y}a\\
Lisela	&	{\hr}raa-n\su{†}	&		&	\cg{(bagu-t)}	&	{\hr}ubaa\su{Δ}	&	\cg{(b-rema-n)}	&	\cg{(deka-t)}\\
Buru	&	{\hr}laa	&	{\hr}dae\su{‡}	&	{\hr}haa-t	&	\cg{(emhalat)}	&	\cg{(eb-rema-n)}	&	\cg{(deka-t)}\\
Ambelau	&	{\hr}haa	&		&	{\hr}e-hea	&	{\hr}buaa	&	\cg{(emrahae)}	&	{\hr}u\tbb{l}an\\	\midrule
Kayeli	&	{\hr}laa-ni	&	{\hr}lae	&	{\hr}lei	&	{\hr}ubaa \it{(met.)}	&	{\hr}hat-lau	&	{\hr}ulan\\
		\lspbottomrule
	\end{tabular}
		\begin{tablenotes}\tnsize
			\item[†]	{Lisela \it{raa-n} ‘sail’ has irregular *l > \it{r}.}
			\item[‡]	{Buru \it{dae} ‘landwards’ is not regular from *daya
								and is probably a loan from an unidentified source.
								The expected reflex *[raa] does not occur.}
			\item[Δ] {Lisela \it{ubaa} `crocodile' is a loan from Kayeli,
								where many crocodiles are present in the swamps
								of the slow-moving lower Wae Apo river.
								(Lisela only has the fast moving Wae Nibe river,
								so not as many crocodiles are present in the
								Lisela-speaking area.)}
		\end{tablenotes}
	\end{threeparttable}
%\end{adjustbox}
\end{table}

\trf{08-09} shows the lexically specific split of *j > \it{y, l}, {\0}.
Two additional examples of *j > \it{y} in Soboyo are *maja >
\it{moyoŋ} ‘dry’, and *bajaq ‘tell, inform; ask,
inquire’ > \it{faya} ‘say (\tsc{1sg}/PL)’,
\it{baya} ‘say (\tsc{2sg}/PL, \tsc{3sg}/PL)’


\begin{table}\tcp{\tsc{Taliabu} and nearby languages *j /V{\gap}V}{08-09}
\begin{adjustbox}{width=\textwidth}
	\begin{threeparttable}\st{2pt}
	\begin{tabular}{llllllll}\lsptoprule
*gloss	&	gall, bile	&	charcoal	&	name	&	sniff	&	ySib\su{‡}	&	how m.?	&	spike trap\\
PMP	&	*qapə\tbr{j}u	&	*qa\tbr{j}əŋ	&	*ŋa\tbr{j}an	&	*ha\tbr{j}ək	&	*hua\tbr{j}i	&	*pi\tbr{j}a	&	*su\tbr{j}a\\	\midrule
Saluan	&	{\hr}pouʔ	&		&	\cg{(saŋgo)}	&	{\hr}ook	&	\cg{(utus)}	&	{\hr}hipian\su{Δ}	&	\\
Balantak	&	(s)opo\tbr{y}uʔ	&	{\hr}ala\tbr{y}oŋ	&	{\hr}ŋaan	&	{\hr}ook	&	\cg{(utus)}	&	{\hr}pi{\tl}pii	&	\\
Banggai	&	\cg{(sopot)}	&	{\hr}u\tbr{y}oŋ	&	\cg{(sambu)}	&	{\hr}o\tbr{y}oki\su{†}	&	{\hr}u\tbb{l}ik	&	{\hr}noian\su{Δ}	&	\\%\hdashline
\noalign{\vskip-1.5ex}
\multicolumn{8}{@{}c@{}}{\dashedline{133mm}} \\
\noalign{\vskip-0.5ex}
Soboyo	&	{\hr}n-o\tbr{y}o-in	&	{\hr}a\tbr{y}oŋ	&	{\hr}ŋaa-in	&	{\hr}ha\tbr{y}oʔ	&	{\hr}u\tbb{l}iʔ, wa\tbb{l}i	&	{\hr}hi\tbb{l}a	&	{\hr}su\tbb{l}a-in\\
Taliabu	&	{\hr}n-o\tbr{y}u-ɲ	&	{\hr}a\tbr{y}oŋ	&	{\hr}ŋaa-ɲ	&	{\hr}ha\tbr{y}ok	&	{\hr}wa\tbb{l}i	&	{\hr}hi\tbb{l}a	&	{\hr}so\tbb{l}a-ɲ\\
Kadai	&	{\hr}n-o\tbr{y}u	&	{\hr}a\tbr{y}oŋ	&	{\hr}ŋaa	&	{\hr}ha\tbr{y}o	&	{\hr}u\tbb{l}i	&		&	\\	\midrule
Mangole	&	{\hr}n-peu	&		&	{\hr}ŋaa	&	\cg{(suun)}	&	\cg{(fuki)}	&	{\hr}pi\tbb{l}a	&	\\
Sanana	&	{\hr}m-peu	&		&	{\hr}naa	&	\cg{(muhi)}	&	\cg{(fuk)}	&	{\hr}bet pi\tbb{l}a	&	{\hr}sua\\
Lisela	&	{\hr}peu-n	&		&	{\hr}ŋaa-n	&		&	{\hr}uwai	&	{\hr}pi\tbb{l}a-ni	&	{\hr}sura-n\su{ǁ}\\
Buru	&	{\hr}peu-n	&	\cg{(ime-n)}	&	{\hr}ŋaa-n	&	{\hr}mae-k	&	{\hr}wai	&	{\hr}pi\tbb{l}a	&	{\hr}sura-n\su{ǁ}\\
Ambelau	&	{\hr}feo-ni	&		&	{\hr}nea-mo	&		&	{\hr}βai-o, eβai-no	&	{\hr}fi\tbb{l}a, fia	&	\\	\midrule
Kayeli	&		&		&	{\hr}naa-ni	&		&	{\hr}wai-n (<Buru)	&	{\hr}hila	&	\\
		\lspbottomrule
	\end{tabular}
		\begin{tablenotes}\tnsize
			\item[†]	{Banggai \it{oyoki} means ‘kiss’.}
			\item[‡]	{In most languages this etymon in this region means ‘younger sibling (same sex)’.
								Taliabu \it{wali} is given for ‘sibling-in-law’,
								Soboyo \it{wali} is given as ‘sister-in-law’.}
			\item[Δ]	{\citet[132]{Mead2003a} posits *i-pija-n ‘when’ as the source
								of Saluan \it{hipian} ‘when’
								and Banggai \it{noian} ‘when (in the future)’.
								\citet[151]{Bergh1953} also gives Banggai \it{iaa} ‘where’}
			\item[{ǁ}]	{If Buru
								and Lisela \it{sura-n} are retentions of *suja they show irregular *j > \it{r}.
								Some words here may be ultimately connected with Sanskrit \it{śūla}
								[ʃuːla] ‘spear, lance’ (see \citealp[568]{Hoogervorst2016}).}
		\end{tablenotes}
	\end{threeparttable}
\end{adjustbox}
\end{table}

Note in \trf{08-09} that lexically specific *j > \it{l} in *pija
is shared with \tsc{Sula-Buru},
but also \tsc{Ambon-Seram} (see \trf{07-05}).
*j > \it{l} in *huaji is shared with Banggai (and with \tsc{Ambon-Seram})
but not with \tsc{Sula-Buru}.
The lexically specific loss of *j > {\0} in *ŋajan is shared
with \tsc{Saluan-Banggai} languages (Andio also has \it{ŋaan}
‘name’ [\citealp[76]{Mead2003b}]) and \tsc{Sula-Buru} (but not with \tsc{Ambon-Seram}).
Note that the underlying sound correspondence of *j > \it{y}
in *qapəju, *qajəŋ, and *hajək is considered to be an important sound change
that defines \tsc{Celebic} languages \citep{Mead2003a}.

Reflexes of *qasawa ‘spouse’ have *awa > \it{oa} in \tsc{Taliabu},
which is shared with \tsc{Saluan-Banggai} \citep[72f, 83]{Mead2003b},
word initially and in the other instance of *awa we have
been able to identify *w = \it{w}.
Examples are give in \trf{08-10}.

\begin{table}\tcp{\tsc{Taliabu} and nearby languages *w}{08-10}
\begin{adjustbox}{width=\textwidth}
	\begin{threeparttable}\st{2pt}
	\begin{tabular}{lllllll}\lsptoprule
*gloss	&	eight	&	vine, rope	&	water	&	mangrove&	spider	&	spouse\\
				&				&							&				&	root		&					&	\\
PMP	&	*\tbr{w}alu	&	*\tbr{w}aRəj	&	*\tbr{w}ahiyəR	&	*\tbr{w}akat\su{†}	&	*la\tbr{w}aq	&	*qasa\tbr{w}a\su{‡}\\	\midrule
Saluan	&	{\hr}\tbr{u}aluʔ	&		&	{\hr}ue, uee	&	{\hr}\tbr{b}akat	&	{\hr}l\tbb{o}aʔ	&	{\hr}os\tbb{o}a\\
Balantak	&	{\hr}\tbr{w}aluʔ	&		&	{\hr}\tbr{w}eer	&	{\hr}\tbr{w}akat	&	{\hr}l\tbb{o}aʔ	&	\cg{(samba-samba)}\\
Banggai	&	{\hh}alu	&	{\hh}alos	&	\cg{(paisu)}	&	\hp{*w}akat	&	{\hr}l\tbb{o}a	&	{\hr}os\tbb{o}aan\\%\hdashline
\noalign{\vskip-1.5ex}
\multicolumn{7}{@{}c@{}}{\dashedline{132mm}} \\
\noalign{\vskip-0.5ex}
Soboyo	&	{\hr}\tbr{w}alu	&	{\hr}\tbr{w}aho	&	{\hr}\tbr{w}ayo	&	{\hr}\tbr{w}aka-ic	&	{\hr}takala\tbr{w}a-ic	&	{\hr}nʃ\tbb{o}a\\
Taliabu	&	{\hr}\tbr{w}alu	&	{\hr}\tbr{w}aho	&		&		&		&	‹tjěhini, tjēhini›\\
Kadai	&	{\hr}\tbr{w}alu	&	{\hr}\tbr{w}aho	&	{\hr}\tbr{w}ayo	&		&	{\hr}takala\tbr{w}a	&	{\hr}nʃ\tbb{o}a\\	\midrule
Mangole	&	\cg{(gata-ua)}	&	\cg{(meu)}	&	{\hr}\tbr{w}ai	&	{\hr}\tbr{w}aka	&	{\hr}la\tbr{w}a	&	{\hr}sau\\
Sanana	&	\cg{(gata-hua)}	&	{\hr}\tbr{w}aʔi	&	{\hr}\tbr{w}ai	&	{\hr}\tbr{w}aka	&	{\hr}la\tbr{w}a	&	{\hr}sau\\
Lisela	&	\cg{(t-rua)}	&	{\hr}\tbr{w}ahe-t	&	{\hr}\tbr{w}ae	&	{\hr}\tbr{w}akat	&		&	{\hr}sa\tbr{w}a-ni\\
Buru	&	\cg{(et-rua)}	&	{\hr}\tbr{w}ahe-t 	&	{\hr}\tbr{w}ae	&	{\hr}\tbr{w}akat	&	\cg{(hakat)}	&	{\hr}em-sa\tbr{w}a-n\\
Ambelau	&	{\hr}\tbr{β}alo, \tbr{w}alo	&	{\hr}\tbr{β}ahe-re	&	{\hr}\tbr{β}ae	&	\cg{(βati)}	&		&	\\	\midrule
Kayeli &	{\hr}\tbr{w}aloʔ	&	{\hr}\tbr{w}ale-t	&	{\hr}\tbr{w}ael-e	&		&	{\hr}lola\tbr{w}a	&	{\hr}sa\tbr{w}a-ne\\
		\lspbottomrule
	\end{tabular}
		\begin{tablenotes}\tnsize
			\item[†]	{The Saluan,
								Balantak, Banggai, Mangole, and Sanana words are given as just ‘root’.
								Soboyo \it{waka-ic} is glossed ‘species of mangrove, \it{Bruguiera gymnorhiza Lam.}’.
								Lisela and Buru \it{wakat} is mangrove.
								Some of these words may be from *wakaR ‘root (generic)’.}
			\item[‡]	{Banggai \it{osoaan} ‘marry’,
								Soboyo and Kadai \it{nʃoa} ‘other half’,
								Mangole and Sanana \it{sau} ‘brother-in-law’,
								Taliabu ‹tjěhini, tjēhini› ‘husband’ \citep[6]{Stokhof1980b},
								Buru \it{emsawan} ‘spouse of child' (defaults to son-in-law.
								Daughter-in-law \it{fina msawan} is usually marked for gender).}
		\end{tablenotes}
	\end{threeparttable}
\end{adjustbox}
\end{table}

We have shown that \tsc{Taliabu} shares three of the sound changes
that define \tsc{Celebic} languages \citep{Mead2003a}: *-ay >
\it{e}, *-aw > \it{o}, *j > \it{y}.
It was also shown that *ə > \it{o,} associated with \tsc{Eastern
Celebic}, is also reflected in \tsc{Taliabu}.
The reduction of medial consonant clusters identified for \tsc{Celebic}
languages \citep{Mead2003a} is also widespread across the
Wallacean subgroups \citep{Blust1993}.
Issues associated with medial consonant clusters
are discussed further in \srf{13-02-01}.
\citet{Mead2003a} also identified PMP *h > {\0}
as a shared innovation of \tsc{Celebic} languages.
While \citet{Blust1981} argues that PMP *h = \it{h} in lexically
specific cases in Soboyo,
the normal reflex is *h > {\0}
and there are other problems/issues with the proposal of *h =
\it{h} which are discussed in greater detail in \srf{13-01-04}.

\tsc{Taliabu} and \tsc{Sula-Buru} share *R > \it{h},
with Mangole and Sanana *R > {\0} as a secondary development,
except that *R > ʔ intervocalically in some dialects of Sanana
in lexically specific ways.\footnote{
\cite[33, 43]{Collins1981} says,
“In the Fagudu and Falahu dialects some intervocalic *R's are reflected as [ʔ]”.
And “Some examples of *R > \it{ʔ} are: *DuRi > \it{hoʔi} ‘thorn’;
*uRat > \it{uʔa} ‘vein’; *biRaq > \it{ka-fiʔa} ‘Alocasia’; *ma-iRaq
> \it{miʔa} ‘red’ and *paRi > \it{paʔi} ‘stingray’.
But in the same dialect we note: *baqeRu > \it{feu} ‘new’; *SaDiRi
> \it{hii} ‘post’; *ñaRa > \it{nai} ‘brother (woman speaking)’.”}
  This can be sequenced as *R > **h > {\0},
(\it{ʔ}), if one assumes they all subgroup together.
\citet{Blust1981} notes that Soboyo \it{h} is from multiple sources,
as from PMP *R, *h, *p, *d.
Yet reflexes of these historical consonants
are kept distinct in \tsc{Sula-Buru} languages.
\trf{08-11} provides additional examples of reflexes of *R.
See also reflexes of *layaR
and *Raya in \trf{08-08},
along with those of *wahiyəR
and *waRəj in \trf{08-10}.

\begin{table}\tcp{\tsc{Taliabu} and nearby languages *R}{08-11}
\begin{adjustbox}{width=\textwidth}
	\begin{threeparttable}\st{2pt}
	\begin{tabular}{lllllll}\lsptoprule
*gloss	&	vein	&	new	&	shoulder\su{†}	&	brother (f.s.)	&	conch	&	male\\
PMP	&	*u\tbr{R}at	&	*baqə\tbr{R}u	&	*qaba\tbr{R}a	&	*ña\tbr{R}a	&	*tabu\tbr{R}i	&	*ma\tbr{R}uqanay\\	\midrule
Saluan	&	{\hr}uat	&	{\hr}buʔou	&	{\hr}oaa	&	\cg{(utus)}	&		&	{\hr}moʔane\\
Balantak	&	{\hr}urat	&	{\hr}uʔuru	&		&	\cg{(utus)}	&		&	{\hr}moroone\\
Banggai	&	\cg{(neet)}	&	\cg{(soodo)}	&	\cg{(beŋe)}	&		&	\cg{(popuukon)}	&	{\hr}malane\\
\noalign{\vskip-1.5ex}
\multicolumn{7}{@{}c@{}}{\dashedline{136mm}} \\
\noalign{\vskip-0.5ex}
Soboyo	&	{\hr}u\tbr{h}a	&	{\hr}fo\tbr{h}u	&	{\hr}fa\tbr{h}a	&	{\hr}na\tbr{h}a	&	{\hr}tafu\tbr{h}i	&	{\hr}mene\\
Taliabu	&	{\hr}n-u\tbr{h}a-c	&	{\hr}fo\tbr{h}u	&	{\hr}fa\tbr{h}a	&	{\hr}na\tbr{h}a	&	{\hr}tafu\tbr{h}i	&	{\hr}mini, molain\\
Kadai	&	{\hr}n-u\tbr{h}a	&	{\hr}fo\tbr{h}u	&	{\hr}fa\tbr{h}a	&	{\hr}na\tbr{h}a	&	{\hr}tafu\tbr{h}i	&	{\hr}mene\\	\midrule
Mangole	&		&	{\hr}f\und{eu}	&		&	{\hr}n\und{ai}	&		&	{\hr}mana\\
Sanana	&	{\hr}\und{ua}	&	{\hr}f\und{eu}	&		&	{\hr}n\und{ai}	&	{\hr}taf\und{oi}	&	{\hr}m\und{aa}na, ma\tbr{ʔ}ana\\
Lisela	&	{\hr}u\tbr{h}a-t	&	{\hr}ɸe\tbr{h}u	&	{\hr}ɸa\tbr{h}a-n	&	{\hr}na\tbr{h}a-t	&		&	{\hr}mana (<Mangole)\\
Buru	&	{\hr}u\tbr{h}a-n	&	{\hr}fe\tbr{h}u-t	&	{\hr}fa\tbr{h}a-n	&	{\hr}na\tbr{h}a-t	&	\cg{(foli)} (< loan)	&	{\hr}em\tbr{h}ana\\
Ambelau	&	{\hr}u\tbr{h}a-re	&	{\hr}bi\tbr{h}u	&	{\hr}ba\tbr{h}a-mu	&	{\hr}na\tbr{h}a-ni	&	{\hr}erb\und{ui}	&	{\hr}e\tbr{h}mana \it{(met.)}\\ \midrule
Kayeli	&	{\hr}ula-t	&	{\hr}belo	&		&	{\hr}nala-ni	&	{\hr}tahuri	&	{\hr}emlana\\
		\lspbottomrule
	\end{tabular}
		\begin{tablenotes}\tnsize
			\item[†] {Soboyo, Taliabu, Kadai \it{faha} ‘arm’ vs.
								\it{lima} ‘hand’; Buru \it{faha-n},
								Lisela \it{ɸaha-n} ‘hand, arm’;
								Ambelau \it{baha-mu} ‘arm’.}
		\end{tablenotes}
	\end{threeparttable}
\end{adjustbox}
\end{table}

The final innovation not yet discussed for \tsc{Celebic} is PMP *d > \it{r}.
In \trf{08-03}, it was shown that in \tsc{Taliabu} PMP *d > \it{h,
d}, and at first glance this does not appear to
share this latter innovation with \tsc{Celebic}.
The regular reflex for \tsc{Taliabu} is *d > \it{h},
while a few reflexes with additional complications do not have
*d > \it{h}: *dahun > Soboyo \it{doiɲ}; \it{kayu ndoiɲ} ‘tree
leaf’, Taliabu \it{n-doŋ}; *ma-diŋdiŋ ‘cold’ > Taliabu \it{bara-tindiŋ} ‘shiver’.
\trf{08-12} shows examples of *d > \it{h}.

Within the region, Saluan
and Sanana also have *d > \it{h},
Ambelau splits *d > \it{l,
h} and *R > \it{h},
and Manusela (\tsc{Ambon-Seram}) has *d/*z > \it{h.} These reflexes
show that *d > \it{h} -- probably through intermediate **r -- is
not an uncommon regional sound change.
Given that the main reflex of PMP *d in \tsc{Taliabu} is \it{h},
then this can be seen as a secondary development through **r
and would be an additional sound change shared with \tsc{Celebic}.
Note that  \tsc{Sula-Buru} keeps reflexes of *d
and *R distinct, even in Sanana where *d > \it{h}, *R > {\0}.

\begin{table}
	\centering\tcp{\tsc{Taliabu} and nearby languages *d}{08-12}
	\st{5pt}\begin{tabular}{llllll}\lsptoprule
*gloss	&	seize	&	two	&	land(ward)	&	stand	&	behind\\
PMP	&	*\tbr{d}akəp	&	*\tbr{d}uha	&	*\tbr{d}aya	&	*kə\tbr{d}əŋ	&	*m-u\tbr{d}əhi\\ \midrule
Saluan	&		&	{\hr}o-\tbr{h}uaʔ	&		&	\cg{(tinʤo)}	&	\cg{(i hiku)}\\
Balantak	&	-\tbb{r}akop	&	{\hr}\tbb{r}uaʔ	&		&	{\hr}-ke\tbb{r}er	&	\\
Banggai	&		&	{\hr}\tbb{l}ua	&		&		&	\\
\dshmidrule{6}
Soboyo	&	{\hr}\tbr{h}ako	&	{\hr}\tbr{h}owo	&	{\hr}\tbr{h}oyo	&	{\hr}ko\tbr{h}o-in	&	{\hr}sa-mu\tbr{h}i\\
Taliabu	&	{\hr}\tbr{h}ako	&	{\hr}\tbr{h}owo	&	{\hr}\tbr{h}oyo	&	{\hr}ko\tbr{h}oŋ	&	{\hr}sa-mu\tbr{h}i \it{‘later’}\\
Kadai	&	{\hr}\tbr{h}ako ‘carry’	&	{\hr}\tbr{h}owo	&		&		&	\\ \midrule
Mangole	&		&	{\hr}ga-ʔuu	&	{\hr}\tbb{l}ai	&	{\hr}ke\tbb{l}i, ge\tbb{l}i	&	\\
Sanana	&	{\hr}\tbr{h}ak kot	&	{\hr}ga-\tbr{h}u	&	{\hr}\tbr{h}ai	&	{\hr}ge\tbr{h}i	&	\\
Lisela	&	\cg{(pogo)}	&	{\hr}\tbb{r}ua	&		&	{\hr}ke\tbb{r}e	&	{\hr}mu\tbb{r}i\\
Buru	&	{\hr}ef-\tbb{r}ake	&	{\hr}\tbb{r}ua	&	\cg{(dae)}	&	{\hr}ke\tbb{r}e-k	&	{\hr}mo\tbb{r}i, mo\tbb{r}i-n\\
Ambelau	&		&	{\hr}\tbb{l}ua	&		&	\cg{(en-βati)}	&	{\hr}mu\tbb{l}i\\ \midrule
Kayeli 	&		&	{\hr}\tbb{l}ua	&	{\hr}\tbb{l}ae	&	\cg{(sabele)}	&	{\hr}bel-mo\tbb{l}iʔ\\
		\lspbottomrule
	\end{tabular}
\end{table}

In a chart of sound correspondences such as \trf{08-03},
one can get the false impression that \tsc{Taliabu}
and Ambelau share PMP *d > \it{h}.
However, the words in which \tsc{Taliabu} *d > \it{h} are different
from the rare Ambelau words in which *d > \it{h}.
For example, PMP *m-udəhi ‘behind’ > Soboyo \it{sa-muhi} ‘behind’
and Taliabu \it{sa-muhi} ‘later’,
but Ambelau \it{muli} ‘behind’.
PMP *daya ‘land(ward)’ > Soboyo,
Taliabu \it{hoyo} ‘landward’, but Ambelau \it{lae} ‘land, landward’.
PMP *duha ‘two’ > Soboyo,
Taliabu, Kadai \it{howo} ‘two’, but Ambelau \it{lua} ‘two’.
Ambelau *d > \it{h} is found in PMP *daRaq ‘blood’ > \it{haha-ni}
‘blood’, possibly due to regressive assimilation (Soboyo \it{poli} Taliabu,
Kadai \it{polo} ‘blood’ is unrelated).\footnote{
		The one other
		possibility in our data is the word for ‘thousand’.
		Keeping in mind that PMP *d > \it{r} in Buru
		and Lisela \it{rara-n} ‘thousand’,
		Ambelau \it{haha-ni} ‘thousand’, and Kayeli \it{lala-ne} ‘thousand’
		seem to suggest an earlier form of **dada(n) ‘thousand’. We have
		not yet found cognates for **dada(n) ‘thousand’ in any other
		language.} 

From the data in \trf{08-11}
and \trf{08-12}, one can propose a merger in \tsc{Taliabu} of
*R/*d > \it{h}, that did not occur in \tsc{Sula-Buru} (*R > \it{h}, *d = **d).

